
我们将通过描述
\[
\mathfrak g_{\ell,k,\beta}
\bigl(\operatorname{Sym}(Y_1,\ldots,Y_\ell)\bigr)
(\mathbf x)\in\overline H[1],
\qquad
\mathbf x\in B_k(\overline H[1]),
\]
的像来定义映射族
$\mathfrak g_{\ell,k,\beta}$。
以后为简便起见，
省略 $\operatorname{Sym}(\ )$，
将其写作
\[
\mathfrak g^\circ_{\ell,k,\beta}
(Y_1,\ldots,Y_\ell)(\mathbf x).
\]
现在说明如何修改 $\mathfrak g$
以得到后文命题 7.4.85 中的
$\mathfrak g^\circ$。

考虑一棵连通平面树\liftIndex{I290}
$\Gamma=(|\Gamma|,Y(\,\cdot\,),\beta(\,\cdot\,))$，
再赋予如下附加数据：
\begin{enumerate}
\item[\textup{(7.4.47.1)}]\label{cond:7-4-47-1}
将顶点集 $C^0(\Gamma)$
分成不交并
\[
C^0(\Gamma)
=C^0_{\mathrm{ext}}(\Gamma)
\cup C^0_M(\Gamma)
\cup C^0_Y(\Gamma).
\]
三个集合中的顶点分别称为
外顶点、$M$ 型顶点与 $Y$ 型顶点。

\item[\textup{(7.4.47.2)}]\label{cond:7-4-47-2}
每个外顶点恰有一条相邻边，
称为外边；其他边称为内边。

\item[\textup{(7.4.47.3)}]\label{cond:7-4-47-3}
外顶点恰有 $k+1$ 个，
依次编号为 $0,1,\ldots,k$。
这一编号遵循平面逆时针定向
所诱导的循环次序。

\item[\textup{(7.4.47.4)}]\label{cond:7-4-47-4}
对每个 $M$ 型顶点 $v$，
指定 $\beta(v)\in G$。
若 $\beta(v)=\beta_0=(0,0)$，
则要求 $v$ 至少邻接三条边。

\item[\textup{(7.4.47.5)}]\label{cond:7-4-47-5}
$Y$ 型顶点恰有 $\ell$ 个。
对每个这样的顶点 $v$，
指定
$Y(v)\in\{Y_1,\ldots,Y_\ell\}$；
映射 $v\mapsto Y(v)$ 给出双射
\[
C^0_Y(\Gamma)
\cong\{Y_1,\ldots,Y_\ell\}.
\]

\item[\textup{(7.4.47.6)}]\label{cond:7-4-47-6}
若 $v$ 是 $Y$ 型顶点，
且 $Y(v)=Y_a$，
则将 $v$ 的相邻边数
记为 $k_a+1$。
\end{enumerate}
置
$\mathbf Y=(Y_1,\ldots,Y_\ell)$。

\begin{definition}{7.4.48}\label{def:7-4-48}
以 $G(k,\mathbf Y)$
表示满足条件 (7.4.47)
的所有 $\Gamma$ 所成的集合。
\end{definition}

当 $\ell=0$，即
$\mathbf Y=\varnothing$ 时，
$G(k,\varnothing)$
与第~\ref{subsec:5-4-4} 小节
所用的树集 $G^+_{k+1}$ 相同。

现在对每个
$\Gamma\in G(k,\mathbf Y)$，
按 $\Gamma$ 的内顶点数归纳定义
两个同态
\[
\mathfrak f(\Gamma):
B_k(\overline H[1])
\longrightarrow\overline C[1],
\qquad
\mathfrak g(\Gamma):
B_k(\overline H[1])
\longrightarrow\overline H[1],
\]
其次数分别为
\[
\deg\mathfrak f(\Gamma)
=\sum_a(\deg Y_a-1),
\qquad
\deg\mathfrak g(\Gamma)
=1+\sum_a(\deg Y_a-1).
\]
这里内顶点按定义就是
$M$ 型或 $Y$ 型顶点。

如前，固定嵌入
$\overline H\hookrightarrow\overline C$，
并将其像记为
$\overline{\mathcal H}$。
再取幂等映射
\[
\Pi_H:\overline C\longrightarrow\overline C,
\qquad
\Pi_H^2=\Pi_H,
\]
使
$\operatorname{Im}\Pi_H
=\overline{\mathcal H}$，
从而
\[
\overline C
=\operatorname{Im}\Pi_H
\oplus\ker\Pi_H.
\]
将 $\overline H$ 与
$\overline C$ 的子空间
$\overline{\mathcal H}$ 识别，
并按引理~\ref{lem:5-4-28}
取链同伦
$\mathfrak G:\overline C\to\overline C$，
使
\[
\overline{\mathfrak m}_1\circ\mathfrak G
+\mathfrak G\circ\overline{\mathfrak m}_1
=-(1-\Pi_H).
\]
借此识别，
也把 $\mathfrak f(\Gamma)$
视为定义在 $\overline C$ 上的映射，
规定它在 $\ker\Pi_H$ 上取零。

先看只有一个内顶点 $v$ 的情形。
设 $v$ 邻接 $k+1$ 条边。
若 $v$ 为 $Y$ 型，置
\begin{equation}
\left\{
\begin{aligned}
\mathfrak f(\Gamma)
&=\mathfrak G\circ Y(v):
B_k(\overline H[1])
\longrightarrow\overline C[1],\\
\mathfrak g(\Gamma)
&=\Pi_H\circ Y(v):
B_k(\overline H[1])
\longrightarrow\overline H[1].
\end{aligned}
\right.
\tag{7.4.49.1}\label{eq:7-4-49-1}
\end{equation}
若 $v$ 为 $M$ 型，置
\begin{equation}
\left\{
\begin{aligned}
\mathfrak f(\Gamma)
&=\mathfrak G\circ
\mathfrak m_{k,\beta(v)}:
B_k(\overline H[1])
\longrightarrow\overline C[1],\\
\mathfrak g(\Gamma)
&=\Pi_H\circ
\mathfrak m_{k,\beta(v)}:
B_k(\overline H[1])
\longrightarrow\overline H[1].
\end{aligned}
\right.
\tag{7.4.49.2}\label{eq:7-4-49-2}
\end{equation}
这里
\[
\mathfrak m_k
=\sum_{\beta\in G}
T^{E(\beta)}e^{\mu(\beta)/2}
\mathfrak m_{k,\beta}.
\]

再设 $\Gamma$ 至少有两个内顶点。
令 $v_0$ 为编号 $0$ 的外顶点，
$e_0$ 为其唯一相邻边，
并记
$\partial e_0=\{v_0,v'_0\}$。
从 $\Gamma$ 中移去
$e_0,v_0,v'_0$，
将余集的各连通分支取闭包，
得到 $\Gamma_1,\ldots,\Gamma_n$。
于是
$\Gamma_i\in G(m_i,\mathbf Y_{(i)})$，
且
$m_1+\cdots+m_n=k$，
并有
\begin{equation}
\left\{
\begin{aligned}
\bigcup_{i=1}^{n}\mathbf Y_{(i)}
&=\mathbf Y,
&&\text{若 $v'_0$ 为 $M$ 型},\\
\bigcup_{i=1}^{n}\mathbf Y_{(i)}
\cup\{Y(v'_0)\}
&=\mathbf Y,
&&\text{若 $v'_0$ 为 $Y$ 型}.
\end{aligned}
\right.
\tag{7.4.50}\label{eq:7-4-50}
\end{equation}
式~\eqref{eq:7-4-50}
左端均为不交并。
按平面的逆时针定向，
依次编号
$\Gamma_1,\ldots,\Gamma_n$
（见图~\ref{fig:5-4-2}）。
若 $v'_0$ 为 $Y$ 型，
定义
\begin{equation}
\left\{
\begin{aligned}
\mathfrak f(\Gamma)
&=\mathfrak G\circ Y(v'_0)\circ
\bigl(
\mathfrak f(\Gamma_1)
\widehat\otimes'\cdots\widehat\otimes'
\mathfrak f(\Gamma_n)
\bigr),\\
\mathfrak g(\Gamma)
&=\Pi_H\circ Y(v'_0)\circ
\bigl(
\mathfrak f(\Gamma_1)
\widehat\otimes'\cdots\widehat\otimes'
\mathfrak f(\Gamma_n)
\bigr).
\end{aligned}
\right.
\tag{7.4.51}\label{eq:7-4-51}
\end{equation}
若 $v'_0$ 为 $M$ 型，
则定义
\begin{equation}
\left\{
\begin{aligned}
\mathfrak f(\Gamma)
&=\mathfrak G\circ
\mathfrak m_{n,\beta(v'_0)}
\circ\bigl(
\mathfrak f(\Gamma_1)
\widehat\otimes'\cdots\widehat\otimes'
\mathfrak f(\Gamma_n)
\bigr),\\
\mathfrak g(\Gamma)
&=\Pi_H\circ
\mathfrak m_{n,\beta(v'_0)}
\circ\bigl(
\mathfrak f(\Gamma_1)
\widehat\otimes'\cdots\widehat\otimes'
\mathfrak f(\Gamma_n)
\bigr).
\end{aligned}
\right.
\tag{7.4.52}\label{eq:7-4-52}
\end{equation}
这里 $\widehat\otimes'$
由下式规定：
\begin{equation}
\begin{aligned}
&(A_1\widehat\otimes'\cdots
\widehat\otimes' A_n)
(u_1\otimes\cdots\otimes u_n)\\
&\qquad
=(-1)^\ast
A_1(u_1)\otimes\cdots\otimes A_n(u_n),
\end{aligned}
\tag{7.4.53}\label{eq:7-4-53}
\end{equation}
其中
\[
\ast=\sum_{i<j}
\deg' A_i\,\deg' u_j.
\]
对
$\Gamma\in G(k,\mathbf Y)$，
还置
\[
\beta(\Gamma)
=\sum_{v\in C^0_M(\Gamma)}\beta(v),
\qquad
E(\Gamma)=E(\beta(\Gamma)),
\qquad
\mu(\Gamma)=\mu(\beta(\Gamma)).
\]

式~\eqref{eq:7-4-54}
中出现的符号
$\ast(\Gamma)$
将在后文式 (7.4.68) 指定。
最后定义
\begin{equation}
\left\{
\begin{aligned}
\mathfrak f(\mathbf Y)(\mathbf x)
&=\sum_{\Gamma\in G(k,\mathbf Y)}
(-1)^{\ast(\Gamma)}
T^{E(\Gamma)}e^{\mu(\Gamma)/2}
\mathfrak f(\Gamma)(\mathbf x),\\
\mathfrak g(\mathbf Y)(\mathbf x)
&=\sum_{\Gamma\in G(k,\mathbf Y)}
(-1)^{\ast(\Gamma)}
T^{E(\Gamma)}e^{\mu(\Gamma)/2}
\mathfrak g(\Gamma)(\mathbf x).
\end{aligned}
\right.
\tag{7.4.54}\label{eq:7-4-54}
\end{equation}
当 $\mathbf Y=\varnothing$ 时，
有
\begin{equation}
\mathfrak f(\varnothing)(\mathbf x)
=\mathfrak f(\mathbf x),
\qquad
\mathfrak g(\varnothing)(\mathbf x)
=\mathfrak m(\mathbf x),
\tag{7.4.55}\label{eq:7-4-55}
\end{equation}
右端由定义~\ref{def:5-4-37} 给出。
下一命题表明，
$\mathfrak g$ 除符号外是
滤过 $L_\infty$ 同态。
